The experiments in this report are restricted to the final publishable artifact set preserved under \path{report/data/sources/}. This artifact set combines the repaired CIFAR-10 retraining run from May 18, 2026, the repaired CIFAR-100 sweep-and-retrain run from May 19--20, 2026, and the stabilized ModelNet40 voxel run from April 29, 2026. The goal is to document what was actually built, debugged, and measured, not to mix the strongest-looking numbers from unrelated runs.

The first workload is CIFAR-10, a 10-class image classification task with \(32 \times 32\) RGB inputs \citep{krizhevsky2009cifar}. The second workload is CIFAR-100, which keeps the same input resolution but increases the label space to 100 classes. Both CIFAR loaders use a deterministic stratified validation split from the training set and preserve the official test split for post-training evaluation. The final CIFAR-10 and CIFAR-100 checkpoints were selected by validation accuracy and then recomputed on the held-out test split.

The third workload is a voxelized ModelNet40 export \citep{wu20153dshapenets}. The repository manifest describes 8{,}858 training samples, 985 train-derived validation samples, and 2{,}468 test samples, with one input channel and a spatial resolution of \(32^3\). This experiment uses \texttt{num\_dims = 3}, which activates the generalized 3D path in patch embedding, window partitioning, shifting, and patch merging. A held-out test recomputation is preserved for the final restored checkpoint, but the report remains explicit that the model was selected by validation accuracy.

\begin{table}[t]
  \centering
  \scriptsize
  \caption{Experimental ledger for the preserved runs. Sweep budgets and summed trial wall-clock times come from the committed sweep summaries; restored epochs and final wall-clock times come from the final \texttt{metrics.json} files.}
  \label{tab:experiment-ledger}
  \begin{tabularx}{\linewidth}{@{}lYYYY@{}}
    \toprule
    \textbf{Workload} & \textbf{Splits used} & \textbf{Sweep budget/time} & \textbf{Final training state} & \textbf{Final wall-clock} \\
    \midrule
    CIFAR-10 & 45k train / 5k validation / 10k test & 48 trials, 60 epochs each; 32.83 h summed & 600-epoch budget; stopped at 392; restored epoch 352 & 14{,}866.06 s (4.13 h) \\
    CIFAR-100 & 45k train / 5k validation / 10k test & 24 trials, 50 epochs each; 13.58 h summed & 600-epoch budget; stopped at 350; restored epoch 300 & 12{,}043.98 s (3.35 h) \\
    ModelNet40 & 8{,}858 train / 985 validation / 2{,}468 test & 15 trials, 10 epochs each; 1.08 h summed & 200-epoch budget; stopped at 27; restored epoch 7 & 990.99 s (0.28 h) \\
    \bottomrule
  \end{tabularx}
\end{table}

\begin{table}[t]
  \centering
  \scriptsize
  \caption{Final retraining configurations. The table keeps the main effective knobs visible while the exact JSON and metric files remain preserved in \autoref{app:artifacts}.}
  \label{tab:retrain-configs}
  \begin{tabularx}{\linewidth}{@{}lYYY@{}}
    \toprule
    \textbf{Workload} & \textbf{Geometry and model} & \textbf{Optimization} & \textbf{Regularization} \\
    \midrule
    CIFAR-10 & 2D; patch \(2 \times 2\); window \(4 \times 4\); embed 96; depths \([2,2,18,2]\) & batch 128; lr \(4.300210 \times 10^{-4}\); wd \(1.041082 \times 10^{-2}\); warmup 20 & drop 0.05; drop path 0.1; cutmix 0.5; color jitter \\
    CIFAR-100 & 2D; patch \(2 \times 2\); window \(4 \times 4\); embed 96; depths \([2,2,18,2]\) & batch 128; lr \(2.047168 \times 10^{-4}\); wd \(2.672492 \times 10^{-2}\); warmup 30 & drop path 0.3; mixup 0.2; cutmix 0.5; cutout 16 \\
    ModelNet40 & 3D; patch \(2 \times 2 \times 2\); window \(4 \times 4 \times 4\); embed 72; depths \([2,2,6,2]\) & batch 64; lr \(1.262274 \times 10^{-3}\); wd \(7.049835 \times 10^{-3}\); warmup 3 & drop path 0.1 \\
    \bottomrule
  \end{tabularx}
\end{table}

Hyperparameter search is performed through configuration-driven random search. The preserved CIFAR-10 sweep evaluated 48 trials with validation accuracy as the selection metric and a 60-epoch budget per trial; trials 3 and 33 tied at the best validation accuracy, and the preserved auto-best configuration uses the first tied trial. The fixed CIFAR-100 sweep evaluated 24 trials with a 50-epoch budget per trial. The ModelNet40 sweep evaluated 15 successful trials with a 10-epoch budget per trial. The sweep plots below show the selected trial for each workload.

\begin{figure}[t]
  \centering
  \begin{tikzpicture}
  \begin{axis}[
    width=0.95\linewidth,
    height=0.35\textheight,
    xlabel={Trial index},
    ylabel={Validation accuracy},
    xmin=-0.5,
    xmax=47.5,
    ymin=0.39,
    ymax=0.75,
    xtick={0,8,16,24,32,40,47},
    ymajorgrids=true,
    grid style={draw=black!10},
    tick label style={font=\small},
    label style={font=\small},
    legend style={
      draw=none,
      fill=none,
      font=\footnotesize,
      at={(0.5,1.14)},
      anchor=south,
      legend columns=2,
      /tikz/every even column/.append style={column sep=0.8em}
    },
    legend cell align=left
  ]
    \addplot[
      color=ndsblue,
      line width=1.1pt,
      mark=*,
      mark size=1.7pt,
      mark options={fill=ndsblue}
    ] table[col sep=comma, x=trial, y=val_accuracy] {data/cifar10_sweep.csv};
    \addlegendentry{Trials}

    \addplot[
      only marks,
      mark=star,
      mark size=4pt,
      color=ndsred,
      fill=ndsred
    ] coordinates {(3,0.7344000339508057)};
    \addlegendentry{Best}

    \node[anchor=south west, font=\scriptsize, text=ndsred] at (axis cs:3.7,0.7350) {trial 3};
  \end{axis}
\end{tikzpicture}

  \caption{Validation accuracy across the 48 preserved CIFAR-10 sweep trials. The highlighted star marks trial 3, one of two tied best trials at 0.7344 under the 60-epoch sweep budget. The later repaired retraining run used this selected artifact as its base but lowered drop-path regularization.}
  \label{fig:cifar10-sweep}
\end{figure}

\begin{figure}[t]
  \centering
  \begin{tikzpicture}
  \begin{axis}[
    width=0.95\linewidth,
    height=0.35\textheight,
    xlabel={Trial index},
    ylabel={Validation accuracy},
    xmin=-0.5,
    xmax=23.5,
    ymin=0.42,
    ymax=0.61,
    xtick={0,4,8,12,16,20,23},
    ymajorgrids=true,
    grid style={draw=black!10},
    tick label style={font=\small},
    label style={font=\small},
    legend style={draw=none, fill=none, font=\small, at={(0.02,0.98)}, anchor=north west},
    legend cell align=left
  ]
    \addplot[
      color=ndsblue,
      line width=1.1pt,
      mark=*,
      mark size=1.7pt,
      mark options={fill=ndsblue}
    ] table[col sep=comma, x=trial, y=val_accuracy] {data/cifar100_sweep.csv};
    \addlegendentry{Trials}

    \addplot[
      only marks,
      mark=star,
      mark size=4pt,
      color=ndsred,
      fill=ndsred
    ] coordinates {(11,0.5902000069618225)};
    \addlegendentry{Best}

    \node[anchor=south west, font=\scriptsize, text=ndsred] at (axis cs:11.5,0.593) {trial 11};
  \end{axis}
\end{tikzpicture}

  \caption{Validation accuracy across the 24 fixed CIFAR-100 sweep trials. The highlighted star marks trial 11, which reached 0.5902 under the 50-epoch sweep budget and was selected for full retraining.}
  \label{fig:cifar100-sweep}
\end{figure}

\begin{figure}[t]
  \centering
  \begin{tikzpicture}
  \begin{axis}[
    width=0.95\linewidth,
    height=0.37\textheight,
    xmin=-0.5,
    xmax=14.5,
    xtick={0,2,4,6,8,10,12,14},
    tick label style={font=\small},
    label style={font=\small},
    grid style={draw=black!10},
    ymajorgrids=true,
    xlabel={Trial index},
    ylabel={Validation accuracy},
    ymin=0.70,
    ymax=0.82
  ]
    \addplot[
      color=ndsblue,
      line width=0.85pt,
      mark=*,
      mark size=2.8pt,
      mark options={fill=ndsblue}
    ] table[col sep=comma, x=trial, y=val_accuracy] {data/modelnet40_sweep_success.csv};

    \addplot[
      only marks,
      mark=star,
      mark size=4.6pt,
      color=ndsred,
      fill=ndsred
    ] coordinates {(11,0.8092095851898193)};

    \node[anchor=west, font=\scriptsize, text=ndsred] at (axis cs:11.55,0.813) {trial 11};
  \end{axis}
\end{tikzpicture}

  \caption{Validation accuracy across the 15 successful ModelNet40 sweep trials in the resumed stable run. The highlighted star marks trial 11, which reached 0.8061 under the 10-epoch sweep budget.}
  \label{fig:modelnet-sweep}
\end{figure}

\clearpage

After sweep selection, each task is converted into the final retraining configuration summarized in \autoref{tab:retrain-configs}. CIFAR-100 preserves the selected drop-path setting, while CIFAR-10 and ModelNet40 use the selected artifacts as bases and lower \texttt{drop\_path\_rate} for the long retraining runs. The restored best state, not the nominal epoch budget, determines the final reported checkpoint. \autoref{tab:experiment-ledger} makes this explicit because the wall-clock time and restored epoch are part of the reproducibility evidence, not just implementation details.

The experiments were executed through the repository's CLI, sweep, queue, and tmux-backed workflow mechanisms, which are themselves part of the practical contribution. The high-level command form for the final sweep-plus-retrain jobs is \texttt{make auto-sweep-tmux SWEEP=configs/sweeps/<sweep>.yaml TRAIN\_EPOCHS=<epochs>}; the detached queue wrapper uses \texttt{make benchmark-tmux} for the committed 2D-to-3D benchmark queue. The committed extraction inputs and the published Hugging Face weight repositories are listed in \autoref{app:artifacts}.
